\documentclass[10pt]{article}

\usepackage[utf8]{inputenc}

\usepackage[T1]{fontenc}

\usepackage{amsmath}

\usepackage{amsfonts}

\usepackage{amssymb}

\usepackage[version=4]{mhchem}

\usepackage{stmaryrd}

\usepackage{bbold}

\usepackage{mathrsfs}

\usepackage{graphicx}

\usepackage[export]{adjustbox}

\graphicspath{ {./images/} }



\begin{document}

тных элементов, касающихся фиксированного подмногообразия $P \subset N$, и точка контакта к-рых принадлежит $P$. Если $P$ - точка, то соответствующие Л. м.- это слой сферизации (или проективизации) касательного расслоения $T^{*} N$.



Другой пример Л. м.- это 1-график функции в контактном многообразии 1 -струй функций $J^{1} N$.



См. также Контактная геометрия. $\quad$ С. А. Табачников. ЛЕЖАНДРОВО ОТОБРАЖЕНИЕ - ограничение на вложенное лежандрово многообразие проекции на базу лежандрового расслоения. Ле жанд ро вой экв и вале нтн о стью Л. о. называется контактоморфизм пространств лежандровых расслоений, переводяший лежандрово многообразие первого Л. о. в лежандрово многообразие второго и слои первого расслоения в слои второго. Примеры Л. о.: нормальное отображение, Лежандра преобразование, тангенциальное отображение (см. Проективная двойственность). Л. о. образуют специальный класс отображений многообразий в многообразия на единицу большей размерности и могут иметь особенности, нетипичные в классе всех отображений. ЛЕЖАНДРОВО РАССЛОЕНИЕ - гЛадкое расслоение контактного многообразия, все слои которого являются лежандровыми подмногообразиями. Проекция в $M$ многообразия контактных элементов многообразия $M$ является Л. р. Все Л. р. одной размерности локально изоморфны. В $\mathbb{R}^{2 n+1}$ (с координатами $x_{1}, \ldots, x_{n}, y_{1}, \ldots, y_{n}, z$ ) с контактной формой $\lambda=d z+\sum x_{i} d y_{i}$ Л. р. задается формулой $(x, y, z) \mapsto(y, z)$. Композиция вложения многообразия в качестве лежандрова и проекции Л.р. на базу называются лежандровым отображением, а его образ - Фр о нтом.



ЛЕМАНА - КеЛЛЕНА CПEKTPAЛЬНОЕ П. П. Табачников. ЛЕНИЕ - см. Спектральное представление.



ЛЕННАРД-ДЖОНСА ПОТЕНЦИАЛ ВЗАИМОДЕЙСТВИЯ - взаимодействия потенциал вида



\[

\Phi=4 \varepsilon\left[\left(\frac{r_{0}}{r}\right)^{12}-\left(\frac{r_{0}}{r}\right)^{6}\right]

\]



Другое название - потен ц и ал $6-12$. Слагаемое $\left(r^{-6}\right)$, описывающее притяжение частиц, обосновано квантовомеханически (дисперсные силы притяжения). Слагаемое $\left(r^{-12}\right)$, описывающее отталкивание на малых расстояниях, не имеет надежного обоснования, но является хорошим приближением (отражает непроницаемость частиц друг для друга). Л.-Д. п. в. имеет минимум, равный $-\varepsilon$ в точке $r=r_{0} \sqrt[6]{2}$



Потенциал взаимодействия принадлежит классу Л.Д. п.в., если он ограничен снизу и существуют такие $0<a_{1}<a_{2}<\infty, A_{1}, A_{2}>0$ и $\lambda_{1}, \lambda_{2}>v(v-$ размерность пространства), что



\[

\begin{gathered}

\Phi(x)>A_{1}|x|^{-\lambda_{1}}, \quad|x|<a_{1}, \\

\Phi(x)>A_{2}|x|^{-\lambda 2}, \quad|x|>a_{2}, \quad x \in \mathbb{R}^{v} .

\end{gathered}

\]



Потенциалы из этого класса сверхустойчивы и регулярны снизу. Если $\Phi(q)<0$ при $|q|>a_{2}$, то Л.-Д. п. в. являются быстро убывающими и регулярными. Использование Л.-Д. п. в. в различных задачах классич. статистич. механики и кинетич. теории газов и жидкостей определяется тем, что он реалистично описывает межмолекулярные взаимодействия (взаимодействия Ван-дер-Ваальса).



Jum.: [1] Lennard - J ones J. E., «Proc. Roy. Soc.», 1924, A 106, p. 463; [2] Бале ску Р., Равновесная и неравновесная статистическая механика, пер. с англ., М., 1978, т. 1, § 6.1; [3] Рюэль Д., Статистическая механика. Строгие результаты, пер. с англ., М., 1971, §3.2.10.



П. В. Малышев. ЛЕПТОКВАРК - сверхтяжелый ( 10 ный или скалярный бозон, возникающий в великого объединения моделях и имеющий вершины взаимодействия одновременно с кварками и лептонами. Л. являются триплетами по цвету и обладают дробным электрич. зарядом. Отличают также д и к в а р к и (секстеты по цвету), взаимодействующие с двумя кварками. Квантовые числа векторных Л. относительно $S U(3) \times S U(2) \times U(1)$ таковы: $(3,2,-5 / 3)$, $(3,2,7 / 3),(3,1,-2 / 3)$. Напр., лагранжиан взаимодействия Л. $X$ с зарядом 4/3 с фермионами имеет вид:



\[

\mathscr{L}=X_{1}^{\mu}\left(A \bar{d}_{R}^{i} \gamma_{\mu} e_{R}^{+}+B \varepsilon_{i j k} \bar{u}_{L_{j}}^{c} \gamma_{\mu} u_{L_{k}}+C \bar{d}_{L}^{i} \gamma_{\mu} l_{L}^{+}\right)

\]



где $A, B, C$ - нек-рые константы, индекс $c$ означает зарядовое сопряжение, индексы $L, R$ - левую и правую спиральность, $i-$ цвет. Л. переносят взаимодействие с несохранением фермионного числа и приводят к pacnaду протона.



М. Е. Шапошников.



ЛЕПТОН - элементарная частица, принимающая участие лишь в слабых и электромагнитных взаимодействиях, но не участвующая в сильных взаимодействиях. К настоящему времени обнаружено три семейства Л., в каждое из к-рых входят электрически заряженный Л. (электрон, мюон и $\tau$-лептон) и соответствующее нейтрино, не несущее электрич. заряда $\left(v_{e}, v_{\mu}, v_{\tau}\right)$. Массы заряженных Л. разных семейств существенно отличаются друг от друга:



\[

m_{\tau}: m_{\mu}: m_{\mathrm{e}} \approx 4000: 200: 1,

\]



тогда как нейтрино либо безмассовы, либо имеют очень маленькую массу. Во всех известных процессах с участием Л. сохраняются по отдельности три квантовых числа: лептонные заряды $L_{\mathrm{e}}, L_{\mu}, L_{\tau}$ (напр., электрон и электронное нейтрино несут единичный электронный лептонный заряд $L_{\mathrm{e}}=1$, мюон и мюонное нейтрино обладают единичным мюонным зарядом, а $\tau$-лептон и $\tau$-нейтрино имеют единичный тауонный заряд). Все свойства Л. находятся в отличном согласии с предсказаниями стандартной модели электрослабых взаимодействий. Однако проблема существования трех семейств Л., во многом похожих друг на друга, но сильно отличающихся по массе, остается в настоящее время нерешенной.



Лит.: [1] Оку н ь Л. Б., Лептоны и кварки, М., 1981; [2] К о мм и н с Ю., Буксбаум Ф., Слабое взаимодействие лептонов и кварков, пер. с англ., М., $1987 . \quad$ А. Ю. Кгнатьев.



ЛЕПТОННЫЙ ТOK - см. Ток.



ЛЕС - см. Боголюбова - Парасюка R-операция.



ЛЕСНАЯ ФОРМУЛА - см. Боголюбова - Парасюка R-операция.



ЛЕСТНИЧНОЕ ПРИБЛИХЕНИЕ - сМ. Бете-СоЛпитера уравнение.



ЛИ АЛTEБPA - векторное пространство, на котором введена операция коммутирования; а именно, любым двум элементам $x, y$ ставится в соответствие третий элемент $[x, y]$, называемый коммутатором, и при этом выполнены аксиомы:



\begin{enumerate}

  \item $[\alpha x+\beta y, z]=\alpha[x, z]+\beta[y, z]$, где $\alpha, \beta$ - скаляры;



  \item $[x, y]=-[y, x]$ (антиком мутати в ность);



\end{enumerate}



\includegraphics[max width=\textwidth, center]{2023_07_08_9900ef8fdab4181ababfg-1}

Якоби).



Если векторное пространство комплексное, то и Л.а. называется комплексной, а если вешественное, то в е щ е с т в е н н о й.



Простейший пример Л. а.- алгебра векторов в трехмерном пространстве с операцией векторного произведения. Другой пример - Л.а., строящаяся по любой ассоциативной алгебре, заменой операции умножения на операцию $[x, y]=x y-y x$.



Главная причина, по к-рой важны Л. а., состоит в том, что по каждой группе Ли однозначно может быть построена алгебра Ли (см. Ли группа, Классические группы и алгебры $\boldsymbol{\Pi u})$.



\includegraphics[max width=\textwidth, center]{2023_07_08_9900ef8fdab4181ababfg-1(1)}

множество, состоящее из всевозможных линейных комбина-





\end{document}